\section*{Abstract}

$\xipar$ ist der einzige fundamentale Parameter der T0-Theorie. Die vollständige Ableitungskette lautet: $\xipar \to \{m_e, m_\mu, \ldots\} \to G \to \lP$, mit $L_0 = \xipar \cdot \lP$ als Sub-Planck-Skala der Raumzeitgranulation. Die maximale kosmologische Skala $R_H$ folgt aus $E_H = \Ezero \cdot \xipar^{41/4}$. Das Verhältnis $R_H/L_0 = 6{,}49 \times 10^{64}$ folgt vollständig aus $\xipar = 4/30000$ allein.

Die Frage nach einer universellen Maximalgröße des Universums ist in der FFGFT jedoch nicht sinnvoll gestellt: \textbf{jede physikalische Skala besitzt ihre eigene systemspezifische Obergrenze} $R_\text{Torus}(m) = \hbar/mc$. Es gibt keine einzige kosmische Obergrenze, die für alle Systeme gilt — $R_H$ ist die Obergrenze des kosmologischen Sektors, nicht des Universums als Ganzes.

\section{Ableitungskette: Von \texorpdfstring{$\xi$}{xi} zu beiden Skalgrenzen}

\subsection{Der einzige fundamentale Parameter}

Die T0-Theorie hat einen einzigen fundamentalen Parameter (Dokument 013):
\begin{equation}
	\boxed{\xipar = \frac{4}{3} \times 10^{-4}}
\end{equation}
Alle physikalischen Größen leiten sich davon ab. Die vollständige Kette lautet:
\begin{equation}
	\xipar \;\to\; \{m_e, m_\mu, m_\tau, \ldots\} \;\to\; G \;\to\; \lP \;\to\; c
\end{equation}
mit $L_0 = \xipar \cdot \lP$ als fundamentale Sub-Planck-Skala der Raumzeitgranulation.

\subsection{Herleitung der Teilchenmassen aus \texorpdfstring{$\xi$}{xi}}

Die Leptonmassen folgen aus $\xipar$ über den geometrischen Quantenzahlenfaktor $f(n,l,j)$ und den T0-Skalierungsfaktor $S_{T0} = 1\,\text{MeV}/c^2$:
\begin{equation}
	m_\ell = \frac{f(n,l,j)^2}{\xipar^2} \cdot S_{T0}
\end{equation}
Das geometrische Mittel der Leptonmassen liefert die charakteristische Energieskala:
\begin{equation}
	\Ezero = \sqrt{m_e \cdot m_\mu} = 7{,}398\,\text{MeV}
\end{equation}
woraus die Feinstrukturkonstante ohne freie Parameter folgt:
\begin{equation}
	\alpha = \xipar \cdot \frac{\Ezero^2}{(1\,\text{MeV})^2} = \frac{1}{137{,}036} \quad \checkmark
\end{equation}

\subsection{Gravitationskonstante und Planck-Länge}

Die Gravitationskonstante folgt aus $\xipar$ und $m_e$ in natürlichen Einheiten (Dokument 013):
\begin{equation}
	G = \frac{\xipar^2}{4\,m_e} \qquad\Rightarrow\qquad \lP = \sqrt{G} = \frac{\xipar}{2\sqrt{m_e}}
\end{equation}
In natürlichen Einheiten mit $m_e = 0{,}511\,\text{MeV}$: $\lP \approx 9{,}33 \times 10^{-5}$ (nat. Einheiten), was nach SI-Umrechnung ergibt:
\begin{equation}
	\lP = 1{,}616 \times 10^{-35}\,\text{m}
\end{equation}

\subsection{Minimale Längenskala}

\begin{keyresult}[Minimale Längenskala aus \texorpdfstring{$\xi$}{xi}]
\begin{equation}
	\boxed{L_0 = \xipar \cdot \lP = 2{,}155 \times 10^{-39}\,\text{m}}
\end{equation}
Die gesamte Ableitungskette $\xipar \to m_e \to G \to \lP \to L_0$ hängt nur von $\xipar$ als einziger fundamentaler Eingabe ab. SI-Zahlenwerte entstehen durch die unvermeidliche Kalibrierung an menschliche Messeinheiten — physikalisch ist die Theorie parameterlos.
\end{keyresult}

\section{Die maximale Skala: Hubble-Radius aus \texorpdfstring{$\xi$}{xi}}

\subsection{Hubble-Energie aus \texorpdfstring{$\xi$}{xi}}

Die Hubble-Energie folgt direkt aus $\xipar$ (Dokument 026):
\begin{equation}
	E_H = \Ezero \cdot \xipar^{41/4}
\end{equation}
mit $\Ezero = \sqrt{m_e \cdot m_\mu} = 7{,}398\,\text{MeV}$ und $\xipar = 4/30000$.

Numerische Auswertung:
\begin{align}
	\xipar^{41/4} &= (1{,}333 \times 10^{-4})^{41/4} = 1{,}908 \times 10^{-40} \\
	E_H &= 7{,}398\,\text{MeV} \times 1{,}908 \times 10^{-40} = 1{,}412 \times 10^{-33}\,\text{eV}
\end{align}

\begin{keyresult}[Hubble-Konstante aus \texorpdfstring{$\xi$}{xi}]
\begin{equation}
	\boxed{H_0^{\text{T0}} = \frac{E_H}{\hbar} = \frac{1{,}412 \times 10^{-33}\,\text{eV}}{6{,}582 \times 10^{-16}\,\text{eV\,s}} \approx 66{,}2\,\text{km/s/Mpc}}
\end{equation}
Abweichung vom Planck-Wert ($67{,}4\,\text{km/s/Mpc}$): $-1{,}9\%$.
\end{keyresult}

\subsection{Hubble-Radius als maximale Skala}

\begin{foundation}[Abgrenzung: Hubble-Radius vs.\ Teilchenhorizont]
In der Standardkosmologie (expandierendes $\Lambda$CDM-Universum) unterscheidet man zwei kosmische Horizonte: den \textit{Hubble-Radius} $R_H = c/H_0 \approx 1{,}4 \times 10^{26}$\,m und den deutlich größeren \textit{Teilchenhorizont} (beobachtbares Universum, $\approx 4{,}4 \times 10^{26}$\,m $\approx 46{,}5$\,Mrd.\,Lj), der die gesamte kausale Vergangenheit seit dem Urknall umfasst. Der Teilchenhorizont ist nur in einem expandierenden Universum mit Urknall sinnvoll definiert.

In der T0-Theorie gibt es keine Expansion und keinen Urknall. Der Hubble-Radius $R_H$ ist hier der einzig relevante geometrische Horizont: jenseits von $R_H$ ergibt die fraktale Wegverlängerung $z \to \infty$, d.\,h.\ keine Information ist mehr empfangbar. Ein Teilchenhorizont im $\Lambda$CDM-Sinne existiert in T0 konzeptuell nicht. Numerische Vergleiche mit Standardwerten sind daher stets auf $R_H$, nicht auf den Teilchenhorizont zu beziehen.
\end{foundation}

Der Hubble-Radius folgt direkt aus der Hubble-Energie:
\begin{equation}
	R_H = \frac{\hbar c}{E_H} = \frac{1{,}973 \times 10^{-7}\,\text{eV\,m}}{1{,}412 \times 10^{-33}\,\text{eV}}
\end{equation}

\begin{keyresult}[Hubble-Radius aus \texorpdfstring{$\xi$}{xi}]
\begin{equation}
	\boxed{R_H = 1{,}398 \times 10^{26}\,\text{m}}
\end{equation}
Vollständig aus $\xipar$ abgeleitet, ohne freie Parameter.
\end{keyresult}

\subsection{Hubble-Radius und Gesamtmasse}

Die Gesamtmasse des beobachtbaren Universums, die dem Hubble-Radius entspricht:
\begin{equation}
	M_U = \frac{R_H \cdot c^2}{2G} = 9{,}43 \times 10^{52}\,\text{kg} = 4{,}7 \times 10^{22}\,M_\odot
\end{equation}

\begin{keyresult}[Hubble-Radius aus \texorpdfstring{$\xi$}{xi}]
\begin{equation}
	\boxed{R_H = 1{,}398 \times 10^{26}\,\text{m}}
\end{equation}
Vollständig aus $\xipar$ abgeleitet, ohne freie Parameter.
\end{keyresult}

\section{Die vollständige Skalenhierarchie}

\subsection{Skalenhierarchie}

\begin{table}[h]
\centering
\begin{tabular}{lll}
\toprule
\textbf{Skala} & \textbf{Wert} & \textbf{Bedeutung} \\
\midrule
$L_0 = \xipar \cdot \lP$ & $2{,}155 \times 10^{-39}\,\text{m}$ & Sub-Planck-Skala \\
$\lP$ & $1{,}616 \times 10^{-35}\,\text{m}$ & Planck-Länge \\
$\lambda_e = \hbar/m_e c$ & $3{,}86 \times 10^{-13}\,\text{m}$ & Elektron-Compton-Wellenlänge \\
$R_H = \hbar c / E_H$ & $1{,}398 \times 10^{26}\,\text{m}$ & Hubble-Horizont \\
\bottomrule
\end{tabular}
\caption{Vollständige Skalenhierarchie aus $\xipar = 4/30000$}
\end{table}

\subsection{Das Verhältnis der Skalen}

\begin{equation}
	\frac{R_H}{L_0} = \frac{1{,}398 \times 10^{26}\,\text{m}}{2{,}155 \times 10^{-39}\,\text{m}}
\end{equation}

\begin{keyresult}[Skalenverhältnis aus \texorpdfstring{$\xi$}{xi}]
\begin{equation}
	\boxed{\frac{R_H}{L_0} = 6{,}49 \times 10^{64}}
\end{equation}
57 Größenordnungen zwischen minimaler und maximaler Skala — beide aus $\xipar = 4/30000$.
\end{keyresult}

\subsection{Die vollständige Ableitungskette}

\begin{keyresult}[Vollständige Kette von \texorpdfstring{$\xi$}{xi} zu beiden Grenzen]
\begin{align}
	&\text{Lagrangian} \;\Rightarrow\; \xipar = 4/30000 \notag \\
	&\quad\Downarrow \notag \\
	&G = \xipar^2/(4m_e) \;\Rightarrow\; \lP = \sqrt{\hbar G/c^3} \notag \\
	&\quad\Downarrow\qquad\qquad\qquad\quad\Downarrow \notag \\
	&L_0 = \xipar \cdot \lP \qquad\qquad\qquad E_H = \Ezero \cdot \xipar^{41/4} \notag \\
	&= 2{,}155 \times 10^{-39}\,\text{m} \qquad\quad\Downarrow \notag \\
	&\qquad\qquad\qquad R_H = \hbar c/E_H \notag \\
	&\qquad\qquad\qquad = 1{,}398 \times 10^{26}\,\text{m} \notag \\
	&\quad\Downarrow \notag \\
	&\boxed{R_H/L_0 = 6{,}49 \times 10^{64} \quad \text{(aus } \xipar \text{ allein)}}
\end{align}
\end{keyresult}

\section{Physikalische Bedeutung}

\subsection{Physikalische Bedeutung von \texorpdfstring{$R_H$}{R\_H}}

In einem statischen Universum ohne Expansion ist der Hubble-Radius der geometrische Horizont, jenseits dessen Photonen im $\xipar$-Feld durch akkumulierte fraktale Wegverlängerung vollständig rotverschoben sind ($z \to \infty$) — nicht durch intrinsischen Energieverlust am Photon, sondern durch geometrische Streckung des Ausbreitungswegs. Lichter\-müdung (\textit{tired light}) tritt nicht auf.

\begin{keypoint}[Rotverschiebung ohne Expansion und ohne Frequenzabhängigkeit]
Die Rotverschiebung entsteht durch fraktale Summation der Wegverlängerung von Photonen im T0-Vakuumfeld: $z \approx \xipar \cdot \ln(d/\lP)$. Da die Wegverlängerung eine Eigenschaft der Geometrie selbst ist — nicht des Photons — nehmen alle Frequenzen exakt dieselbe Streckung. Es gibt keine Frequenzabhängigkeit. Der Effekt ist dem expandierenden Universum beobachtungsmäßig äquivalent, aber mechanistisch fundamental verschieden. Ausführlich hergeleitet in Dokument 026 (Geometrische Kosmologie).
\end{keypoint}

\subsection{Skalenrelativität des Torus: fraktale Verschachtelung}

In der T0-Theorie ist der 4D-Torus keine kosmische Behälterstruktur, in der sich Teilchen befinden. Jedes physikalische System trägt vielmehr \textbf{seinen eigenen charakteristischen Torus}, dessen Skala durch die Compton-Wellenlänge $\bar\lambda = \hbar/mc$ des jeweiligen Systems bestimmt wird:

\begin{equation}
	R_{\text{Torus}}(m) = \frac{\hbar}{mc}
\end{equation}

\begin{table}[h]
\centering
\begin{tabular}{lll}
\toprule
\textbf{System} & \textbf{Charakteristische Skala} & \textbf{Verhältnis zu $L_0$} \\
\midrule
Elektron     & $3{,}86 \times 10^{-13}$\,m & $\sim 10^{26}$ \\
Proton       & $2{,}10 \times 10^{-16}$\,m & $\sim 10^{23}$ \\
Wasserstoff  & $5{,}29 \times 10^{-11}$\,m & $\sim 10^{28}$ \\
Galaxie      & $\sim 10^{21}$\,m           & $\sim 10^{60}$ \\
Universum    & $R_H = 1{,}40 \times 10^{26}$\,m & $6{,}49 \times 10^{64}$ \\
\bottomrule
\end{tabular}
\caption{Systemspezifische Torus-Skalen — alle mit $L_0$ als gemeinsamer Untergrenze}
\end{table}

\begin{keypoint}[Fraktale Verschachtelung ohne starre Hierarchie]
Diese Strukturen sind nicht wie russische Puppen ineinandergesteckt — der Elektronen-Torus sitzt nicht \textit{im} Galaxien-Torus. Jede Skala ist die geometrische Eigenstruktur des jeweiligen Systems, konstituiert durch dessen Masse. Die gemeinsame Untergrenze $L_0$ ist auf allen Ebenen dieselbe; die Obergrenze ist systemspezifisch.

Die fraktale Kopplung zwischen den Skalen ist durch $\xipar$ vermittelt und damit sehr klein. Wechselwirkungen zwischen weit voneinander entfernten Skalen — etwa zwischen Elektronenstruktur und kosmologischem Feld — sind nicht null, aber unterdrückt durch Potenzen von $\xipar \approx 1{,}33 \times 10^{-4}$. Diese schwache skalenübergreifende Kopplung ist physikalisch relevant: sie ist der Ursprung der fraktalen Korrekturfaktoren $K_{\text{frak}}$ in der Berechnung von $g$-Faktoren, Teilchenmassen und $G$ — kleine, aber messbare Beiträge, die im Standardmodell fehlen oder als freie Parameter eingeführt werden müssen.
\end{keypoint}

\subsection{Lokalität, flacher Torus und Endlichkeit der Berechnungen}

Der T0-Torus ist nicht der anschauliche Donut-Torus, der in $\mathbb{R}^3$ eingebettet ist. Er ist der \textbf{flache Identifikations-Torus} $T^n = \mathbb{R}^n / \mathbb{Z}^n$: ein euklidischer Raum, bei dem gegenüberliegende Ränder identifiziert werden. Dieser Torus hat \textbf{intrinsische Krümmung $K = 0$ überall punktweise} — nicht nur im Integral. Lokal ist er exakt identisch mit $\mathbb{R}^n$, ohne jede Näherung.

\begin{keypoint}[Flacher Torus: lokal exakt euklidisch]
Beim eingebetteten Donut-Torus in $\mathbb{R}^3$ variiert die Gaußsche Krümmung: außen positiv ($K > 0$), innen negativ ($K < 0$). Das Integral verschwindet exakt (Gauss-Bonnet: $\int K\,dA = 2\pi\chi = 0$ für $\chi = 0$). Beim flachen Identifikations-Torus hingegen gilt $K = 0$ \emph{überall} — es gibt keine Innen- und Außenkrümmung, die sich erst im Integral aufheben müssen. Berechnungen innerhalb dieses Torus sind exakt dieselben wie in unbegrenztem euklidischen Raum. Keine Krümmungskorrekturen, keine topologischen Terme.
\end{keypoint}

Daraus folgen drei Konsequenzen für Berechnungen in der T0-Theorie:

\textbf{1. Keine Renormierung notwendig.} Alle Feldintegrale laufen über den endlichen lokalen Torus mit Skala $R_\text{Torus}(m) = \hbar/mc$. Da dieser Torus endlich ist und lokal flach, sind UV-Divergenzen strukturell ausgeschlossen. Die obere Integrationsgrenze ist nicht $\infty$, sondern $R_\text{Torus}$; die untere ist $L_0$. Beide sind endlich und aus $\xipar$ abgeleitet.

\textbf{2. Die $\xipar$-Korrektur für unendliches Universum ist nicht vernachlässigbar klein, sondern schlicht irrelevant.} Formal würde ein unendliches Universum eine Korrektur $\sim \xipar \cdot \ln(R_\infty / R_H)$ erzeugen, die mit $R_\infty \to \infty$ logarithmisch divergiert. Diese Divergenz greift aber nicht, weil alle physikalischen Observablen Integrale über den \emph{lokalen} Torus des jeweiligen Systems sind. Beiträge von Skalen $\gg R_\text{Torus}(m)$ koppeln nicht an lokale Physik — sie sind durch $\xipar$-Potenzen unterdrückt und fallen außerhalb des Integrationsbereichs. Die Frage, ob das Universum endlich oder unendlich ist, hat keinen Einfluss auf lokale Berechnungen.

\textbf{3. Krümmungskorrekturen sind absolut vernachlässigbar.} Die relative Krümmungskorrektur auf der kleinsten physikalischen Skala $L_0$ innerhalb eines Torus der Skala $R_\text{Torus}$ beträgt $(L_0 / R_\text{Torus})^2$:
\begin{align}
	\text{Elektron:} \quad & \left(\frac{L_0}{\bar\lambda_e}\right)^2 = \left(\frac{2{,}155\times 10^{-39}}{3{,}86\times 10^{-13}}\right)^2 \approx 3{,}1 \times 10^{-53} \notag \\
	\text{Proton:} \quad & \left(\frac{L_0}{\bar\lambda_p}\right)^2 \approx 1{,}1 \times 10^{-46} \notag
\end{align}
Diese Werte sind für alle praktischen Berechnungen exakt null.

\begin{keypoint}[Ein Feld — Massen als Knoten — emergente Kohärenzlänge]
Hinter diesen drei Konsequenzen steht eine einheitliche ontologische Struktur: Es gibt \textbf{ein einziges $\xipar$-Vakuumfeld}. Es gibt keine separate Hierarchie von Feldern für verschiedene Skalen.

Massen sind \textbf{topologische Knoten} in diesem Feld — lokale Strukturen, keine externen Objekte. Jeder Knoten der Masse $m$ trägt durch die Dualität $\tilde{T} = 1/m$ eine charakteristische Zeitskala, und daraus emergiert unmittelbar seine räumliche Kohärenzlänge $\hbar/mc$. Der Torus-Radius ist nicht von außen zugewiesen, sondern folgt aus der Knotenstruktur selbst.

Die Hierarchie der Torus-Skalen — vom Proton-Knoten ($\sim 10^{-16}$\,m) bis zum kosmologischen Horizont ($\sim 10^{26}$\,m) — ist damit nichts anderes als die \textbf{fraktale Selbstähnlichkeit eines einzigen Feldes}: dasselbe $\xipar$ auf jeder Skala, dieselbe Untergrenze $L_0$, systemspezifische Kohärenzlängen als emergente Eigenschaft der jeweiligen Knotenstruktur. Die skalenübergreifende Kopplung ist nicht null, aber durch $\xipar$-Potenzen unterdrückt — daher rechnet man lokal, ohne das Gesamtfeld zu kennen, und erhält trotzdem exakte Ergebnisse.
\end{keypoint}

 des T0-Feldes über kosmologische Skalen und bedarf noch der vollständigen Ableitung aus der $\xipar$-Feldtheorie (Dokument 026). Die numerische Übereinstimmung mit $H_0 = 66{,}2\,\text{km/s/Mpc}$ ($-1{,}9\%$ von Planck) ist ein starkes Konsistenzargument.

\begin{foundation}[Zur Rolle von $R_H$ in Berechnungen]
Der genaue Zahlenwert von $R_H$ spielt für die physikalischen Berechnungen innerhalb der T0-Theorie eine untergeordnete Rolle. Die strukturell bedeutsame Größe ist $L_0 = \xipar \cdot \lP$ — sie folgt rein aus $\xipar$ und ist die eigentliche Grundlage aller Skalenaussagen.

$R_H$ hingegen ist in zweifacher Hinsicht anders geartet: Erstens ist sein Zahlenwert an den gemessenen $H_0$ gekoppelt (bzw.\ hängt vom noch abzuleitenden Exponenten $41/4$ ab). Zweitens ist $R_H$ in T0 nicht primär ein berechnetes Ergebnis, sondern eine \textbf{Definitionsgrenze}: die Skala, jenseits derer $z \to \infty$ gilt und keine Beobachtung mehr möglich ist. Mehr als $R_H$ ist per Konstruktion des Modells nicht zugänglich — unabhängig davon, ob der Exponent $41/4$ exakt hergeleitet oder an $H_0$ gefittet wird.

Das Verhältnis $R_H/L_0 = 6{,}49 \times 10^{64}$ drückt daher keine tiefe theoretische Wahrheit aus, sondern den Abstand zwischen der geometrisch abgeleiteten Untergrenze und dem beobachtungsbasierten Horizont. Die eigentliche Leistung der Theorie liegt in $L_0$ — nicht in $R_H$.
\end{foundation}

\subsection{Kosmologische Konstante aus \texorpdfstring{$\xi$}{xi}}

Die dimensionslose kosmologische Konstante folgt direkt:
\begin{equation}
	\Lambda_{\text{dim}} = \left(\frac{\lP}{R_H}\right)^2 = \left(\frac{1{,}616 \times 10^{-35}}{1{,}398 \times 10^{26}}\right)^2 \approx 1{,}34 \times 10^{-122}
\end{equation}
Der Standardwert beträgt $\Lambda_{\text{dim}} \approx 2{,}87 \times 10^{-122}$ (aus $\Lambda \approx 1{,}1 \times 10^{-52}\,\text{m}^{-2}$, CODATA). Das T0-Ergebnis liegt in derselben Größenordnung (Faktor~2{,}1) — keine Dunkle Energie, sondern eine geometrische Konsequenz der Skalenhierarchie aus $\xipar$. Die vollständige Übereinstimmung erfordert die noch ausstehende Ableitung des Exponenten $41/4$ (Dokument 026).

\section{Zusammenfassung}

\begin{keyresult}[Vollständige Skalenhierarchie aus \texorpdfstring{$\xipar = 4/30000$}{xi = 4/30000}]
\begin{align}
	\text{Untergrenze:} \quad & L_0 = \xipar \cdot \lP = 2{,}155 \times 10^{-39}\,\text{m} \notag \\
	\text{Planck:} \quad & \lP = 1{,}616 \times 10^{-35}\,\text{m} \notag \\
	\text{Elektron:} \quad & \lambda_e = \hbar/m_e c = 3{,}86 \times 10^{-13}\,\text{m} \notag \\
	\text{Hubble:} \quad & R_H = \hbar c/E_H = 1{,}398 \times 10^{26}\,\text{m} \notag
\end{align}
\begin{equation}
	\boxed{\text{Alle Skalen aus } \xipar = 4/30000 — \text{kein freier Parameter}}
\end{equation}
Das Universum ist geometrisch selbstkonsistent geschlossen — nicht durch einen Rand, sondern durch die Topologie des 4D-Torus mit $L_0$ und $R_H$ als natürlichen Grenzen derselben Geometrie.

\begin{keyresult}[Kernaussage: Gibt es eine Maximalgröße des Universums?]
Die Frage ist in der FFGFT nicht sinnvoll als globale Frage gestellt. Jede physikalische Skala — vom Proton bis zum kosmologischen Horizont — besitzt ihre \textbf{eigene systemspezifische Obergrenze} $R_\text{Torus}(m) = \hbar/mc$, emergent aus der Knotenstruktur der jeweiligen Masse im einzigen $\xipar$-Feld. $R_H$ ist die Obergrenze des kosmologischen Sektors. Die universelle Untergrenze $L_0$ gilt für alle Skalen gleichermaßen — eine universelle Obergrenze existiert nicht, weil es keine skalenunabhängige Physik gibt, für die sie relevant wäre.
\end{keyresult}
\end{keyresult}


\begin{foundation}[Epistemologische Rahmung: Modellgrenzen statt Existenzgrenzen]
$L_0$ und $R_H$ sind keine physikalischen Wände und keine Aussagen über die Gesamtgröße oder Gesamtstruktur des Universums. Sie sind \textbf{Grenzen der Beschreibbarkeit} innerhalb der T0-Theorie: unterhalb von $L_0$ divergiert die Raumzeitgranulation, jenseits von $R_H$ divergiert die Rotverschiebung ($z \to \infty$). Was außerhalb dieser Grenzen existiert oder nicht existiert, ist keine Frage der Theorie.

Der entscheidende Vorzug gegenüber anderen Ansätzen liegt nicht in dieser epistemologischen Bescheidenheit allein, sondern darin, dass T0 \textbf{keinerlei Regularisierungstricks benötigt}: keine Renormierungsverfahren zur Beseitigung von UV-Divergenzen, keine künstlichen Cutoffs, keine Counterterms, keine Einführung von Zusatzfeldern zur Kontrolle des Unendlichen. $L_0$ entsteht nicht als nachträgliche Korrektur, sondern als direkte geometrische Konsequenz von $\xipar$. Die Theorie ist von vornherein endlich — nicht weil Unendlichkeiten weggerechnet wurden, sondern weil der Parameter $\xipar$ sie strukturell ausschließt.
\end{foundation}

\begin{foundation}[Hinweis: Verbindung zum SI-System]
Die absoluten Zahlenwerte in diesem Dokument — $L_0$, $R_H$, $M_U$ in SI-Einheiten — setzen eine Verbindung zum SI-Einheitensystem voraus. Diese ist nicht zirkulär, sondern epistemologisch notwendig: Aus einem dimensionslosen $\xipar$ allein lassen sich keine absoluten Meter- oder Kilogramm-Werte ableiten. Die Verbindung erfolgt über $m_e$, $\lP$ oder $\alpha$ als äquivalente Einstiegspunkte. Ausführlich behandelt in Dokument 180 (Abschnitt: Verbindung zum SI-System) sowie in Dokument 056 und 101.

Dokumente 180, 181 und 182 bilden eine Einheit und sollten nicht getrennt gelesen werden.
\end{foundation}

\vspace{1cm}
\noindent\textit{Der Exponent $41/4$ bedarf noch der vollständigen Ableitung aus der $\xipar$-Feldtheorie (Dokument 026).}
